% from ega4-2-fr.tex
\oldpage[IV-2]{144}

\begin{remarks}[6.5.5]
\EGAReaderTarget{ega4.c04.item.6-5-5}{IV.2.s6.5.5-fr}
\medskip\noindent
\begin{enumerate}
\item[(i)] Soient $A$, $B$ deux anneaux locaux noethériens,
$\varphi:A\to B$ un homomorphisme local tel que $B$ soit un $A$-module
plat. Soit $k$ le corps résiduel de $A$, et supposons que les deux anneaux
$A$ et $B\otimes_Ak$ vérifient la propriété $(R_i)$ (5.8.2); alors il ne
s'ensuit pas nécessairement que $B$ vérifie $(R_i)$, même dans le cas
particulier où $i=0$ ou $i=1$ et où $B$ est le complété $\widehat A$ de
$A$. Nagata a en effet donné un exemple où $A$ est normal (donc vérifie
$(R_1)$) mais où $\widehat A$ n'est même pas réduit (donc ne vérifie pas
$(R_0)$) [30]. On ne peut appliquer à ce cas la proposition (6.5.3)
parce que les fibres $f^{-1}(y)$ ne vérifient pas nécessairement la
propriété $(R_i)$ aux points distincts du point fermé de $Y=\operatorname{Spec}(A)$.
Nous montrerons toutefois plus loin (7.8.3, (v)) que de tels phénomènes
ne se présentent pas pour les anneaux locaux noethériens que l'on rencontre
le plus souvent dans les applications.

\item[(ii)] La propriété pour un préschéma d'être \emph{intègre} ne se
comporte pas du tout comme les propriétés que nous venons d'examiner dans
ce numéro et les précédents: il peut se faire que
$f:X\to Y$ soit un morphisme plat de type fini, que toutes les fibres
$f^{-1}(y)$ soient régulières (et même géométriquement régulières (6.7.6))
et que $Y$ soit intègre sans que $X$ soit même localement intègre.
Par exemple, soit $k$ un corps algébriquement clos de caractéristique $0$,
et soit $A=k[U,V]/(UV-(U+V)^3)$, l'anneau intègre (dont le spectre est
donc une «~cubique à point double~»); $A$ n'est pas intégralement clos,
et si $u$, $v$ sont les images canoniques de $U$, $V$ dans $A$, la clôture
intégrale de $A$ est l'anneau $C=A[t]$, avec
\[
t=(u-v)/(u+v),\qquad t^2=1+u+v,
\]
d'où l'on tire
\[
u=\frac{1}{2}(t^3+t^2-t),\qquad
v=\frac{1}{2}(-t^3+t^2+t),
\]
et par suite $C=k[t]$, isomorphe à l'anneau de polynômes à une
indéterminée sur $k$. Si $\mathfrak m=Au+Av$, idéal maximal de $A$
(correspondant au «~point double~» de la cubique), il y a deux idéaux
maximaux $\mathfrak n_1$, $\mathfrak n_2$ de $C$, engendrés respectivement
par $t-1$ et $t+1$. Soit alors $B$ le sous-anneau du produit $C\times C$
formé des couples de polynômes $(f,g)$ tels que
$f(1)=g(-1)$ et $f(-1)=g(1)$
($\operatorname{Spec}(B)$ est le schéma obtenu en «~recollant~» deux
exemplaires de $\operatorname{Spec}(C)$, le point $\mathfrak n_1$
(resp. $\mathfrak n_2$) de l'un des deux exemplaires étant «~recollé~»
au point $\mathfrak n_2$ (resp. $\mathfrak n_1$) de l'autre; cf. chap.~V,
où cette opération sera discutée de façon générale). Il y a donc deux
idéaux maximaux $\mathfrak r_1$, $\mathfrak r_2$ de $B$ au-dessus de
l'idéal maximal $\mathfrak m$ de $A$. En outre, comme le processus de
«~recollement~» permute à la localisation et à la complétion, on vérifie
aisément que les homomorphismes canoniques
\[
\widehat A_{\mathfrak m}\to\widehat B_{\mathfrak r_1},
\qquad
\widehat A_{\mathfrak m}\to\widehat B_{\mathfrak r_2}
\]
sont bijectifs, et par suite (Bourbaki, \emph{Alg. comm.}, chap.~III,
\S~3, n\textsuperscript{o}~5, prop.~10) $B_{\mathfrak r_1}$ et
$B_{\mathfrak r_2}$ sont des $A_{\mathfrak m}$-modules plats, ayant
d'ailleurs même corps résiduel que $A_{\mathfrak m}$. Pour tout autre idéal
maximal $\mathfrak p$ de $A$, il est immédiat qu'il y a deux idéaux
maximaux $\mathfrak q_1$, $\mathfrak q_2$ de $B$ au-dessus de $\mathfrak p$,
et que les homomorphismes
$A_{\mathfrak p}\to B_{\mathfrak q_1}$ et
$A_{\mathfrak p}\to B_{\mathfrak q_2}$ sont bijectifs. On voit ainsi que
le morphisme $\operatorname{Spec}(B)\to\operatorname{Spec}(A)$ est plat
et fini, et que toutes ses fibres sont géométriquement régulières (il est
même \emph{\'etale}, comme nous le verrons plus tard (17.6.3)); cependant
il est immédiat que $B$ n'est pas intègre.
\end{enumerate}
\end{remarks}

